% Generado por sistema/backend ReporteTex.java -- no editar a mano.
% Datos: Base de datos: 148 importado (titulados) / Base de datos: 53 importado (empleadores)
\textit{Paso 1. Contar con la muestra y determinar su tamaño.} Solo responden quienes trabajan en una organización. Cada rango se reemplaza por su punto medio en meses ($x_i$); ``más de 12 meses'' se toma como 12, por lo que la media es una cota inferior. El tamaño de la muestra es $n = 101$.

\begin{mitabla}{Cálculo de la media y la desviación estándar de tiempo hasta el primer empleo}
\begin{tabular}{lccccc}
\toprule
Respuesta & $x_i$ & $f_i$ & $f_i\,x_i$ & $(x_i-\bar{x})^2$ & $f_i\,(x_i-\bar{x})^2$ \\
\midrule
Ya trabajaba & 0 & 53 & 0.00 & 6.9362 & 367.6177 \\
$<$ 1 mes & 0.5 & 10 & 5.00 & 4.5525 & 45.5252 \\
1–4 meses & 2.5 & 14 & 35.00 & 0.0179 & 0.2501 \\
4–8 meses & 6 & 6 & 36.00 & 11.3322 & 67.9933 \\
8–12 meses & 10 & 13 & 130.00 & 54.2629 & 705.4179 \\
$>$ 12 meses & 12 & 5 & 60.00 & 87.7283 & 438.6413 \\
\midrule
$\Sigma$ & & 101 & 266.00 & & 1625.4455 \\
\bottomrule
\end{tabular}
\end{mitabla}

\textit{Paso 2. Calcular la media muestral.}
\[
\bar{x} = \frac{\sum f_i\,x_i}{n} = \frac{266.00}{101} = 2.6337\ \text{meses}
\]

\textit{Paso 3. Calcular la desviación estándar muestral} (Coss Bu, ecuación 6.5).
\[
S = \sqrt{\frac{\sum f_i\,(x_i-\bar{x})^2}{n-1}} = \sqrt{\frac{1625.4455}{100}} = 4.0317
\]

\textit{Paso 4. Determinar el estadístico teórico.} Con $\alpha = 0.05$ y $n-1 = 100$ grados de libertad, como los grados de libertad son más de 30 la distribución $t$ se aproxima con la normal, igual que en la prueba de los promedios: $t_{100,\,0.025} \approx z_{0.025} = 1.96$.

\textit{Paso 5. Calcular el margen de error.}
\[
\frac{S}{\sqrt{n}}\,t_{n-1,\,\alpha/2} = \frac{4.0317}{\sqrt{101}}\times 1.960 = 0.4012 \times 1.960 = 0.7863
\]

\textit{Paso 6. Construir el intervalo y concluir.}
\[
2.6337 - 0.7863 \;\le\; \mu \;\le\; 2.6337 + 0.7863 \quad\Longrightarrow\quad 1.85 \le \mu \le 3.42
\]

Con un 95\% de confianza, la media de tiempo hasta el primer empleo de todos los titulados está entre 1.85 y 3.42 meses.

